% introduction.tex -- BNR-2026-013

\section{Introduction}
\label{sec:introduction}

Zero-knowledge proof systems underpin the security guarantees of
validity rollups, enabling off-chain computation with on-chain
verification. The choice of proof system has cascading effects on
system architecture: it determines circuit design methodology,
verification cost, proof aggregation strategy, and operational
complexity~\cite{plonk,groth16}.

Basis Network is an enterprise-grade Avalanche L1 (Subnet-EVM)
providing zero-fee permissioned blockchain infrastructure for
industrial traceability in Latin America. The current Validium MVP
employs Groth16~\cite{groth16} via Circom/SnarkJS, with a production
circuit of 274,291 R1CS constraints achieving 12.9-second proving
time and 306K gas on-chain verification~\cite{basis-validium}. While
Groth16 delivers the smallest known SNARK proofs (128 bytes, constant),
its per-circuit trusted setup requirement creates significant
operational overhead: every modification to the circuit necessitates
a new Phase~2 ceremony.

The evolution from application-specific validium to a general-purpose
zkEVM L2---where each enterprise deploys its own chain with full EVM
compatibility---demands a proof system that supports:

\begin{enumerate}
  \item \textbf{Universal setup}: A single Structured Reference String
    (SRS) reusable across all circuit configurations, eliminating
    per-circuit ceremonies.
  \item \textbf{Custom gates}: PLONKish arithmetization with
    arbitrary-degree polynomial constraints, enabling dedicated gates
    for EVM opcodes and ZK-friendly hash functions.
  \item \textbf{EVM-compatible verification}: Proofs verifiable via
    BN254 pairing precompiles (EIP-196, EIP-197) within the 500K gas
    budget of the Basis Network L1.
  \item \textbf{Rust-native implementation}: Alignment with the
    architectural decision (TD-002) to implement the zkEVM prover
    in Rust for performance and memory safety.
  \item \textbf{Recursive proof composition}: Native support for
    proof aggregation to batch multiple enterprise proofs into a
    single L1 verification.
\end{enumerate}

This paper presents a systematic evaluation of five candidate proof
systems against these requirements, culminating in benchmarks
comparing Groth16 and halo2-KZG across arithmetic and hash chain
circuits. We demonstrate that halo2-KZG meets all enterprise targets
and propose a phased migration strategy with dual verification to
preserve backward compatibility.

\subsection{Contributions}

The contributions of this paper are:

\begin{itemize}
  \item A systematic elimination analysis of five ZK proof systems
    (Groth16, PLONK-KZG, halo2-KZG, halo2-IPA, plonky2) against
    enterprise zkEVM constraints.
  \item Rust benchmarks (30 iterations, release mode) comparing
    Groth16 (arkworks) and halo2-KZG proving time, verification time,
    and proof size across 7~circuit configurations.
  \item Quantification of custom gate constraint reduction: 2.4$\times$
    measured for degree-5 S-box gates, 17$\times$ projected for full
    Poseidon hashing.
  \item A dual-verification migration architecture enabling
    incremental transition from Groth16 to PLONK without service
    interruption.
\end{itemize}

\subsection{Paper Organization}

Section~\ref{sec:related} surveys related work across proof systems
and production zkEVM deployments. Section~\ref{sec:system} defines the
system model and enterprise constraints. Section~\ref{sec:methodology}
describes the experimental methodology. Section~\ref{sec:results}
presents benchmark results and the proof system selection analysis.
Section~\ref{sec:discussion} discusses implications, custom gate
design, and migration strategy. Section~\ref{sec:conclusion}
concludes with recommendations for downstream implementation.
